% discussion.tex — Discussion section for Sentinel Lattice paper
\section{Discussion}
\label{sec:discussion}

\subsection{Limitations}

Our evaluation carries several important caveats that bound the
generalizability of the reported results.

\paragraph{Simulation-based evaluation.}
All attack and defense measurements reported in Section~\ref{sec:evaluation}
derive from simulated adversarial interactions, not production deployments.
Real-world attack distributions may differ substantially in sophistication,
frequency, and composition.  We caution against treating simulated detection
rates as production guarantees.

\paragraph{Partial implementation.}
Only L1 (\sentinel{} Core, comprising 704 detection patterns across 53
engines) is fully implemented.  The upper layers---L2 through
\mire{}---exist as formal designs with projected performance
characteristics.  Until these layers are deployed and evaluated under
realistic conditions, performance claims for the complete lattice
remain projections rather than empirical findings.

\paragraph{\pasr{} key management.}
The HMAC key underpinning \pasr{} provenance signatures constitutes a
single point of failure: compromise of this key invalidates all
downstream trust decisions.  We mitigate this risk through hardware
security modules (HSMs) and ephemeral key rotation, but these
mechanisms add operational complexity that may limit adoption in
resource-constrained deployments.

\paragraph{Red team findings.}
Three attack classes exposed fundamental weaknesses in the current
architecture.  The Landauer cost bound enforced by \cafl{}
(COMBO~ALPHA) proved largely decorative: red team exercise ATK-014
achieved 80--90\% attacker success, confirming that the thermodynamic
cost of bit erasure is orthogonal to semantic danger assessment.
Multi-session accumulation (ATK-012) yielded 55--70\% attacker success,
demonstrating that per-session isolation is insufficient against
persistent adversaries who distribute malicious payloads across
interactions.  Defense-as-DoS (ATK-020) achieved 60--75\% success by
exploiting the nine-layer architecture as a cost amplifier---each
layer adding computational overhead that an attacker can deliberately
trigger.

\subsection{Three Irreducible Residual Classes}

Despite the layered defenses of the \sentinel{} Lattice, our analysis
identifies three classes of residual risk that no purely algorithmic
defense can eliminate.  Together, these account for the entirety of
unmitigated attack surface.

\paragraph{Semantic identity (${\sim}35\%$ of residual risk).}
When malicious intent is linguistically indistinguishable from benign
intent, no text-analysis algorithm can reliably discriminate between
them.  A request for ``how to synthesize compound~X'' may originate
from a chemistry student or a threat actor; the surface forms are
identical.  Resolving this ambiguity requires human judgment or
contextual signals---organizational role, access history, declared
purpose---that lie outside the scope of prompt analysis.  \tsa{} and
\gps{} reduce the incidence of such cases but cannot eliminate them.

\paragraph{Model trust chain (${\sim}45\%$ of residual risk).}
If the underlying language model is compromised prior to
deployment---through training data poisoning, weight manipulation, or
supply-chain attacks---Goldwasser and Kim~\cite{goldwasser2022} prove
that detecting such compromise from black-box query access alone is
impossible.  \mire{} provides containment by isolating model outputs
and enforcing output-space invariants, but it cannot verify the
integrity of internal model representations.  This residual class is
the dominant source of unmitigated risk.

\paragraph{Representation gap (${\sim}20\%$ of residual risk).}
Attacks encoded in modalities that the \pasr{} transducer does not
fully analyze---embedded images, audio segments, or video
frames---bypass text-centric detection layers entirely.  The current
architecture processes only textual and structured-data inputs;
multimodal payloads receive only surface-level metadata tagging.
Closing this gap requires extending the transducer pipeline with
modality-specific analyzers, which we discuss in
Section~\ref{sec:future-work}.

\subsection{\pasr{} Known Weaknesses}

Beyond the key-management concern noted above, \pasr{} exhibits three
structural weaknesses that merit explicit acknowledgment.

\paragraph{Cross-source composition.}
When multiple provenance sources combine within a single prompt---e.g.,
a user instruction wrapping a retrieved document containing
tool-generated output---field-level trust assignment defaults to
conservative (lowest-trust) propagation.  This policy is safe but
lossy: legitimate composite inputs may receive artificially depressed
trust scores, increasing false-positive rates.

\paragraph{Two-model divergence.}
The \pasr{} transducer and the primary LLM may interpret identical
token sequences with divergent semantics.  If the transducer parses a
passage as benign while the LLM interprets it as an instruction,
provenance annotations will misattribute the resulting behavior.  This
divergence is difficult to detect without shared internal
representations between the two models.

\paragraph{Provenance is not authorization.}
\pasr{} establishes \emph{who} contributed each segment of input, but
provenance alone does not determine \emph{whether} that contribution
should be permitted.  Authorization requires policy layers---such as
\aas{} and \irm{}---that consume provenance metadata as one signal
among many.  Deployments that treat provenance as a sufficient access
control mechanism will inherit a false sense of security.

\subsection{Future Work}
\label{sec:future-work}

Several directions extend the contributions of this paper.

\paragraph{L2--\mire{} implementation.}
The most immediate priority is a production-grade implementation of the
upper lattice layers (L2 through \mire{}) in Rust, targeting memory
safety and low-latency inference-path integration.  We estimate
approximately 16 engineer-weeks for a feature-complete prototype
suitable for controlled deployment.

\paragraph{Production deployment and real-world evaluation.}
Transitioning from simulated to production evaluation is essential for
validating detection rates, latency overhead, and false-positive
behavior under authentic workloads.  We plan staged rollouts with
instrumented monitoring to capture attack distributions that simulation
cannot reproduce.

\paragraph{Multimodal extension.}
Extending the \pasr{} transducer to analyze image, audio, and video
inputs would address the representation-gap residual class identified
above.  This requires modality-specific feature extractors that produce
provenance-compatible annotations consumable by downstream layers
including \tsa{} and \gps{}.

\paragraph{Formal verification.}
We intend to pursue machine-checked proofs of \tsa{} safety
properties---specifically, that well-typed provenance annotations are
preserved through the full inference pipeline.  Candidate frameworks
include Lean~4 and F* for the core type-theoretic invariants.

\paragraph{Framework integration.}
Practical adoption requires first-class integration with widely used
LLM orchestration frameworks such as LangChain and LlamaIndex.
Middleware adapters that expose \sentinel{} primitives (\asra{},
\tcsa{}, \pasr{}) as composable pipeline stages would lower the
barrier to deployment.

\paragraph{Adaptive adversarial robustness.}
Finally, we plan systematic adversarial robustness testing against
adaptive attackers who specifically target the defense architecture
itself---attackers who have read this paper and can craft inputs
designed to exploit the lattice structure.  The defense-as-DoS
(ATK-020) finding underscores the urgency of this direction.
